\documentclass[11pt]{article}

\usepackage[margin=1.1in]{geometry}
\usepackage{amsmath,amssymb,amsthm}
\usepackage{booktabs,longtable,array}
\usepackage{microtype}
\usepackage{listings}
\usepackage[hidelinks]{hyperref}
\usepackage{url}
\hypersetup{
  pdftitle={Reconstructing x87 FPATAN on Intel Skylake},
  pdfauthor={},
  pdfsubject={Validated retained-width behavioral reconstruction},
  pdfkeywords={x87, FPATAN, Skylake, floating-point, reverse engineering}
}

\lstdefinestyle{pseudocode}{
  language=Python,
  basicstyle=\ttfamily\footnotesize,
  keywordstyle=\bfseries,
  commentstyle=\normalfont\itshape\footnotesize,
  columns=fullflexible,
  keepspaces=true,
  showstringspaces=false,
  breaklines=true,
  frame=lines,
  aboveskip=0.6\baselineskip,
  belowskip=0.6\baselineskip
}
\newcommand{\T}[1]{\operatorname{T}_{#1}}
\newcommand{\N}{\operatorname{RN}_{64}}
\newcommand{\rc}{\mathrm{RC}}
\newcommand{\insn}[1]{\texttt{#1}}
\newcommand{\hex}[1]{\texttt{#1}}
\newtheorem{proposition}{Proposition}
% Generated from authenticated evidence; do not hand-edit.
\newcommand{\CorpusRows}{2{,}783{,}208}
\newcommand{\ProspectiveRows}{624{,}312}
\newcommand{\RowsPerMode}{695{,}802}
\newcommand{\PCtwentyfour}{98{,}908}
\newcommand{\PCfiftythree}{98{,}908}
\newcommand{\PCsixtyfour}{2{,}585{,}392}
\newcommand{\FinalChallengeRows}{363{,}256}


\title{Reconstructing x87 FPATAN on Intel Skylake:\\
A Retained-Width, Bit-Exact Behavioral Program}
\author{}
\date{September 2026}

\begin{document}
\maketitle

\begin{abstract}
We present a fixed numerical reconstruction of the x87 \insn{FPATAN}
instruction on a Skylake Xeon reference environment. The program uses the
physically decoded Pentium constant ROM, exact magnitude ordering and
table-index selection, distinct direct and table-assisted kernels, and
explicit 64- and 67-bit internal rounding operations. Both kernels contain
two interleaved polynomial chains; their square uses a 67-bit operand and
its 64-bit truncated counterpart. The architectural wrapper covers raw80
operand classes, all four rounding modes, \insn{C1}, and masked arithmetic
exception flags under a stated clear-state contract. No input-indexed
correction, fitted exception selector, host arctangent, or algorithm flag
is used. The final program matches all \CorpusRows{} retained observations
from fourteen one-shot campaigns, including \ProspectiveRows{} prospective
observations of the unchanged final graph. We also give a finite exact
certificate proving that lower-tie and odd-tie table indexing are
endpoint-equivalent within this graph. The result is an empirically
validated behavioral reconstruction, not an exhaustive all-input hardware
proof or a recovered physical microcode implementation. Complete programmer's
pseudocode and the literal constants accompany the specification.
\end{abstract}

\section{Scope and contributions}

The x87 \insn{FPATAN} instruction consumes $y$ in \insn{ST(1)} and $x$ in
\insn{ST(0)}, computes a quadrant-sensitive arctangent, and pops the stack
once. It therefore corresponds to the two-argument role of \texttt{atan2},
not simply to a division followed by a one-argument library
arctangent~\cite{intel2}. This paper concerns the function of the raw input
bits evaluated by the reference processor. It does not replace that target
with correctly rounded mathematical $\operatorname{atan2}(y,x)$.

The reconstruction, designated V7 in the accompanying artifacts, consists
of one numerical graph and one masked architectural wrapper. The
contributions are (i) an explicit retained-width specification, (ii) literal
constant provenance separated from behavioral inference, (iii) prospective
and retained verification with output and flag accounting, (iv) an exact
index-representation equivalence result, and (v) a runnable C implementation
and shared, executable reference pseudocode.

The API assumes a valid two-deep x87 stack, all exceptions masked, and clear
initial exception latches. Its outputs are the raw80 numerical result,
\insn{C1}, and the six arithmetic exception flags. It does not model arbitrary
saved-state restoration, unmasked exception delivery, stack faults, FPU
pointers, or undefined condition bits. Generality here means a single
arithmetic program over operand classes, not a per-input correction table;
the finite corpus is not an enumeration of the input space.

\section{Provenance and experimental method}
\label{sec:method}

\subsection{The observed reference}

Captures identify the reference as an Intel Xeon Processor (Skylake, IBRS),
with CPUID leaf-1 signature \hex{00050654} and reported microcode \hex{0x1}.
The Linux capture environment reports the hypervisor feature bit. These
are recorded guest-visible identifiers, not independent attestation of
the underlying host's physical microcode revision. Claims in this paper
refer to this observed environment; transfer to other CPUIDs, microcode
versions, or vendors has not been established for this FPATAN program.

The harness executes \insn{FNINIT}, loads the requested control word, loads
the two raw80 operands, records the pre-instruction status, executes one
\insn{FPATAN}, and records the result and status. Control words and the
stack-pop relation are checked by the capture protocol. There is no warmup
instruction. A hardware observation is identified by the ordered raw
operands, RC, PC, and capture/CPU context. Frozen inputs and predictions
precede fresh observation; durable reservations prevent repetition.
Existing captures are reused for software regression, not relabeled as new
hardware evidence.

\subsection{Literal ROM data}
\label{sec:rom}

Shirriff's physical P5 ROM decode supplies the constant bit patterns, not
coefficients fitted to the observed misses~\cite{shirriff}. The implementation
contains 44 entries: rows 19 and 20 for $\pi$ and $\pi/2$; rows 114--117 and
118--123 for two polynomial coefficient sets; and rows 125--156 for the
32 arctangent table anchors. The V7 table path uses anchors $n=2,\ldots,32$;
the $n=1$ entry is retained as literal data.

For published exponent $E_i$, sign $s_i$, and hexadecimal significand
integer $S_i$, the value used is exactly
\begin{equation}
 C_i=(-1)^{s_i} S_i\,2^{E_i-65535-66}.
 \label{eq:rom}
\end{equation}
All selected flag bits are zero. The significand has 67 value bits; its
string representation avoids an accidental 64-bit conversion. All 44
entries have been checked against both the archived published appendix and
the local transcription. Appendix~\ref{app:constants} reproduces the exact
values in this encoding. This establishes P5 provenance, not a physical
dump of Skylake's ROM.

\subsection{Independent structural context}

A pinned public Goldmont listing and FP-ROM projection provide a separate
structural clue~\cite{goldmont,fprom}. The ten coefficient payloads and 32
table payloads agree after projection to 64 bits. An arctangent-shaped
long-polynomial block at \hex{U70ca--U70e4} and a short block at
\hex{U6d1e--U6d32} use two interleaved chains, distinct sum-operation classes,
and a distinct square operation. The audit establishes 42 payload
projections and operation incidence; it does not establish discarded low
bits, exponent/sign metadata, dynamic architectural dispatch, or Skylake
micro-operation semantics. The numerical widths below are part of the
behavioral reconstruction and its separate validation.

Story and Tang's IA-64 compatibility work documents magnitude ordering,
direct and table-assisted arctangent paths, and a residual of the form
$(v-Bu)/(u+Bv)$~\cite{story}. Its thresholds, formats and corrections are
different. It is historical algorithmic context, not the source of a
Skylake bit-exactness claim or an error bound for this program.

\section{The fixed numerical graph}
\label{sec:graph}

\subsection{Exact values and explicit rounding}

A finite raw80 encoding has sign $s$, exponent field $e_f$, and integer
significand $m$. Its value is
\begin{equation}
 x=(-1)^s m\,2^{\max(e_f,1)-16383-63}.
\end{equation}
Classification precedes decoding: a nonzero exponent with a zero explicit
integer bit is unsupported. Exponent-zero nonzero encodings include both
subnormals and pseudo-denormals, whose values use effective exponent one.

Let $\T{p}(v)$ truncate toward zero to $p$ significant bits, with unbounded
internal exponent range, and let $\N(v)$ round to nearest/even at 64 bits.
Define
\begin{equation}
 M(a,b)=\T{67}(ab),\qquad
 A(a,b)=\N(a+b),\qquad
 W(a,b)=\T{67}(a+b).
 \label{eq:ops}
\end{equation}
Every arithmetic operation without an explicit rounding operator is exact.
These internal operators are fixed: they do not inherit architectural RC
or PC. Final architectural rounding is specified separately.

\subsection{Magnitude ordering and reduction}

For nonzero finite operands let $a=\min(|y|,|x|)$,
$b=\max(|y|,|x|)$, and $r=a/b$. Retain whether $|y|>|x|$.
The first-octant angle $t$ follows three paths:
\begin{enumerate}
\item If $r<2^{-40}$, set $t=\T{67}(r)$.
\item If $2^{-40}\le r\le3/64$, set $z=\T{67}(r)$ and use the direct kernel.
\item Otherwise select
\begin{equation}
 n=\left\lceil32r-\tfrac12\right\rceil,\qquad c=n/32,\qquad 2\le n\le32,
 \label{eq:index}
\end{equation}
and set
\begin{equation}
 z=\T{67}\!\left(\frac{\T{67}(a-cb)}{\T{67}(b+ca)}\right).
 \label{eq:reduce}
\end{equation}
Use the table kernel, then set $t=\T{67}(\mathrm{kernel})+C_{124+n}$.
\end{enumerate}
Equation~\eqref{eq:index} is nearest-cell selection with ties to the lower
cell. Equation~\eqref{eq:reduce} completes the small-integer products before
cutting either sum. Cutting $ca$ or $cb$ first, or replacing this operation
with a separately rounded quotient-based identity, is a different program.

\subsection{Interleaved polynomial evaluation}

Both kernels first compute
\begin{equation}
 u=\N\bigl(z\,\T{64}(z)\bigr),\qquad v=M(u,u).
 \label{eq:square}
\end{equation}
The square is not $\N(z^2)$ in general: one operand is truncated to 64 bits
before the exact product is rounded. The direct kernel uses
\begin{align}
 o&=W\bigl(C_{119},M(v,A(C_{121},M(v,C_{123})))\bigr),\\
 e&=W\bigl(C_{118},M(v,A(C_{120},M(v,C_{122})))\bigr),
\end{align}
while the table kernel uses
\begin{equation}
 e=W(C_{114},M(v,C_{116})),\qquad o=A(C_{115},M(v,C_{117})).
\end{equation}
Both finish with
\begin{equation}
 h=A(M(u,o),e),\qquad q=M(M(z,u),h),\qquad
 \mathrm{kernel}=z+q.
 \label{eq:terminal}
\end{equation}
In exact algebra these are odd polynomials with interleaved coefficient
chains. With the stated cuts, algebraically equivalent schedules need not
be numerically equivalent. In particular, the direct kernel's terminal
$z+q$ remains exact until the next specified stage. The table kernel is
instead cut to 67 bits before its ROM-anchor addition.

\subsection{Quadrant restoration}

If a magnitude swap occurred or $x$ is negative, first replace
$t$ by $\bar t=\T{67}(t)$. Before applying the sign of $y$, use
\begin{equation}
 \theta_+=
 \begin{cases}
  t & |y|\le|x|,\ x>0,\\
  C_{20}-\bar t & |y|>|x|,\ x>0,\\
  C_{20}+\bar t & |y|>|x|,\ x<0,\\
  C_{19}-\bar t & |y|\le|x|,\ x<0.
 \end{cases}
\end{equation}
Set $\theta=\operatorname{sign}(y)\theta_+$. These additions use the literal
ROM constants and exact arithmetic. They are not separate architectural
floating-point instructions with RC-dependent intermediate rounding.

\section{Architectural encoding, classes and flags}
\label{sec:architecture}

For a nonzero retained angle $\theta$, final spacing is
\begin{equation}
 \Delta=2^{\max(\lfloor\log_2|\theta|\rfloor,-16382)-63}.
\end{equation}
Round $\theta$ to a multiple of $\Delta$ using requested RC, then normalize
and pack the raw80 result. RN breaks ties toward even retained significands;
RD, RU and RZ round toward negative infinity, positive infinity and zero.
Negative nonzero values retain their sign when rounded to zero. The
reconstructed precision-control independence means PC24/53/64 are accepted
architectural controls, not switches between numerical kernels.

\insn{C1} is one exactly when final rounding increased the magnitude:
\begin{equation}
 \insn{C1}=[\,|\operatorname{round}_{\rc}(\theta)|>|\theta|\,].
\end{equation}
Every ordinary nonzero finite path sets PE. It also sets UE when the
retained angle is tiny \emph{before} final rounding,
$|\theta|<2^{-16382}$, including directed rounding up to minimum normal.
This is a statement of the reconstructed instruction behavior, not a
substitution of generic mathematical-library exception rules.

The special-class handler implements the quadrant-sensitive zero and
infinity responses and x87 NaN-priority rules~\cite{intel2,intel1}.
Unsupported encodings return \hex{ffff:c000000000000000} and IE.
Signaling NaNs raise IE and are quieted. With two NaNs, a quiet NaN takes
priority over a signaling NaN; within the same class the larger significand
wins, with positive sign breaking an equal-significand tie. Unsupported
encoding handling precedes NaN propagation in this model.

On non-NaN, supported paths, an exponent-zero nonzero operand, including a
pseudo-denormal, sets DE. Exact zero results from zero/infinity paths do
not set PE; nonzero quadrant constants do. ZE and OE remain clear under
this contract. Pre-load exception flags are checked and are zero for the
tested raw80 load sequence. Appendix~\ref{app:pseudocode} gives the complete
ordered class dispatch, including signed-zero cases; it should be used
instead of inferring a wrapper from the finite formulas alone.

\paragraph{A visible RC example.}
For $x=y=1$, RN and RU return \hex{3ffe:c90fdaa22168c235} with C1=1;
RD and RZ return \hex{3ffe:c90fdaa22168c234} with C1=0. PE is set and
the other arithmetic exception flags are clear. All three supported PC
settings give the same corresponding results.

\section{Validation and discriminating evidence}
\label{sec:validation}

Table~\ref{tab:campaigns} partitions all retained observations; its rows are
not additional counts to sum with the total. Each observation includes an
RC/PC setting, so it is not a count of distinct unconditioned operand pairs.
Only D0023, D0024 and D0026 were prospective tests of the unchanged V7
program. The earlier observations became discovery/regression evidence for
that program, even where they had been prospective for a predecessor.
The earlier eleven campaigns are D0001--D0009, D0013 and D0022.

\begin{table}[htbp]
\centering
% Generated; observation counts are not unconditioned pair counts.
\begin{tabular}{@{}lrl@{}}
\toprule
Evidence set & Observations & Role for the final graph \\
\midrule
Earlier eleven campaigns & 2{,}158{,}896 & Discovery/regression \\
D0023: exact midpoints & 194{,}808 & Prospective V7 \\
D0024: denominator low bits & 66{,}248 & Prospective V7 \\
D0026: final-rounding boundaries & 363{,}256 & Prospective V7 \\
\midrule
Total & 2{,}783{,}208 & Fourteen campaigns \\
\bottomrule
\end{tabular}

\caption{Retained evidence, partitioned by campaign. Every row has zero
final-V7 output, C1, exception and pre-load-flag misses.}
\label{tab:campaigns}
\end{table}

The immutable corpus contains \CorpusRows{} observations. RC coverage is
balanced at \RowsPerMode{} observations each for RN, RD, RU and RZ.
PC coverage is \PCtwentyfour{} at PC24, \PCfiftythree{} at PC53 and
\PCsixtyfour{} at PC64. These counts show the actual tested controls; they
do not claim every possible class/RC/PC combination has been exhaustively
enumerated. Some canonical both-special pairs were conservatively held by
the history audit and have only documented/synthetic coverage, not invented
native observations.

The prospective challenges target exact midpoints with varied denominators,
unmasked denominator low bits, signs and octants, and final-rounding
transitions. D0026 contains 2,048 independently checked adjacent external
RN/RD transition brackets, five nearby significands per bracket, sign/swap
expansions and common exponent transports, plus independent full-exponent
controls. Its \FinalChallengeRows{} observations have zero misses, 1,984
underflow observations and 1,404 complete three-PC comparison groups without
value/status differences. Those groups overlap the observation count.

\subsection{A selector control, not an input exception}

A preceding upper-tie program failed at the exact ratio $19/64$. The final
program uses one global lower-tie rule rather than a special case for that
ratio. The later low-bit campaign independently distinguishes several
frozen alternative index constructions. In D0024, upper and even tie rules
each disagree on 1,007 observations in output or C1; the tested CHOP67
reciprocal/CHOP67 multiply/upper-tie construction disagrees on 935. Lower
and odd tie rules have no disagreement. These are exclusions of the tested
constructions, not of every conceivable reciprocal or index implementation.
The remaining lower/odd equivalence has an exact explanation in
Section~\ref{sec:alias}.

\subsection{Implementation and pseudocode checks}

The main C source, GCC 15 build and independent library client replay all
fourteen jobs with no output or flag differences. Later readability changes
preserve the full preprocessed compiler-token stream and all existing code
comments. An optimized Clang executable and a library build with address
and undefined-behavior sanitizers again reproduce the per-job output hashes over
all \CorpusRows{} observations. Build variants and replays are software
checks of the same observations, not new hardware experiments.

The programmer's pseudocode is shared literally between the standalone
reference and Appendix~\ref{app:pseudocode}. Its exact-rational blocks are
executed by a separate documentation verifier without calling the production
numerical functions. That verifier checks all \CorpusRows{} retained
observations and reproduces the delivered C output-stream hashes. Source
and evidence hashes are preserved in the accompanying artifact manifests.

\section{An exact index-representation equivalence}
\label{sec:alias}

\begin{proposition}
Within the fixed graph in Section~\ref{sec:graph} and its architectural
wrapper, nearest table indexing with ties to the lower integer and nearest
indexing with ties to the odd integer produce identical modeled endpoints
and applicable flags for all valid raw80 operand pairs.
\end{proposition}

\begin{proof}
The selectors are identical away from exact midpoints. At an odd lower
index they also agree. At $1/64$ both possible indices dispatch to the same
direct kernel. Thus only even lower indices $n=2,4,\ldots,30$ can matter,
with exact ratio $r=m/64$, $m=2n+1$.

Normalize the larger magnitude as $b=S2^e$, where
$2^{63}\le S<2^{64}$; this includes exact normalization of subnormals.
For either neighboring cell $j\in\{n,n+1\}$, the numerator in
Equation~\eqref{eq:reduce} is exactly $\pm b/64$, hence unaffected by its
67-bit cut. Put $K=2048+jm$. The denominator can be written
\begin{equation}
 d=(KS-\delta)2^{e-11},\qquad 0\le\delta\le511,
\end{equation}
because $K<4096$, $KS<2^{76}$, and retaining 67 bits discards at most nine
integer bits. Therefore its pre-division-cut reduced magnitude satisfies
\begin{equation}
 \frac{32}{K}\ \le\ \frac{32S}{KS-\delta}\ \le\
 \frac{32\,2^{63}}{K\,2^{63}-511}.
 \label{eq:enclosure}
\end{equation}
The upper bound uses the smallest allowed $S$ and largest allowed
$\delta$. The enclosures deliberately include states that need not have
representable external preimages.

Enumerating every C67 result in these thirty signed neighbor enclosures
gives 106 retained reduced states. The accompanying exact-integer/rational
certificate evaluates their short kernels, table additions, all four
magnitude restorations, both signs and all four rounding modes: 3,392
endpoint checks. Within each midpoint group, every state in both neighbor
enclosures has the same output/C1 vector. All such angles are normal and
nonzero, so modeled PE and UE agree; DE depends on the unchanged input
classes. Other operand classes share the same wrapper. This completes the
finite certificate for the stated graph.
\end{proof}

This proposition does not identify a physical selector in the processor.
Nor does it prove that the graph agrees with silicon outside the observed
corpus. It does show why seeking an output separator between these two
representations under the same graph cannot succeed.

\section{Implementation, reproducibility and limits}

The C11 implementation uses GMP for exact integer/rational arithmetic, not
for an arctangent function. The historical filename
\path{fpatan_candidate.c} now contains the fixed program; its executable
has no algorithm-selection flags or environment-controlled corrections.
The optional library wraps the same implementation. It runs on non-x86
hosts without native FPATAN. Architectural RC/PC inputs are part of the
public contract, not configuration patches.

The accompanying directory contains the source, standalone/library build,
\path{PSEUDOCODE.md}, this manuscript, a complete literal-constant table,
and the evidence/build instructions. The corpus records retain input,
prediction, hardware, identity and reservation hashes. The exact index
certificate is reproducible without a solver or new hardware. No private
supplemental ledger, captured-operand correction table or undisclosed
coefficient fitting is needed to execute the numerical program.

The limits are substantive: corpus agreement is not a proof over all raw80
pairs; one observed reference is not cross-generation validation; static
Goldmont structure is not a Skylake decode; the finite index certificate
proves an identity of two model representations, not hardware equivalence;
and the numerical API is not a complete x87 machine-state emulator. No
bound on error relative to mathematical arctangent is claimed here. The
ongoing adversarial design and unexecuted large campaigns are deliberately
not presented as results.

\section{Conclusion}

The verified program provides an explicit, reusable description of observed
FPATAN behavior: ROM-based reduction, two interleaved retained-width
kernels, ordered quadrant restoration and architectural rounding/flags.
Its \CorpusRows{} matching observations include \ProspectiveRows{}
prospective tests of the unchanged final graph. The specification and
shared pseudocode make the rounding decisions auditable without reading a
research chronology. The numerical reconstruction is validated within its
stated scope; physical uniqueness and universal hardware equivalence remain
distinct, unclaimed results.

\begin{thebibliography}{9}
\bibitem{intel2} Intel Corporation.
\emph{Intel 64 and IA-32 Architectures Software Developer's Manual},
Volume 2A, \insn{FPATAN} and Table 3-30.
\href{https://www.intel.com/content/dam/www/public/us/en/documents/manuals/64-ia-32-architectures-software-developer-vol-2a-manual.pdf}{Official instruction-set reference}.

\bibitem{intel1} Intel Corporation.
\emph{Intel 64 and IA-32 Architectures Software Developer's Manual},
Volume 1, revision 090, February 2026, Table 4-8 and rounding control.
\href{https://cdrdv2-public.intel.com/874241/253665-090-sdm-vol-1.pdf}{Official architecture manual}.

\bibitem{shirriff} K. Shirriff.
\emph{Pi in the Pentium: reverse-engineering the constants in its
floating-point unit}, January 2025, ROM appendix.
\href{https://www.righto.com/2025/01/pentium-floating-point-ROM.html}{Original physical reverse-engineering report}.

\bibitem{goldmont} chip-red-pill.
\emph{uCodeDisasm: Goldmont operation listing},
\path{ucode/ucode_glm.txt}, commit \hex{ffc9070}.
\href{https://github.com/chip-red-pill/uCodeDisasm/blob/ffc9070233a6e7a26dbabe723289259f087ee20b/ucode/ucode_glm.txt}{Pinned primary source}.

\bibitem{fprom} P. Borrello and contributors.
\emph{CustomProcessingUnit}, \path{bios/dumps/rom.txt},
commit \hex{4237524}.
\href{https://github.com/pietroborrello/CustomProcessingUnit/blob/4237524fe7545c66e42dd986113f220662c06f6a/bios/dumps/rom.txt}{Pinned FP-ROM projection}.

\bibitem{story} S. Story and P. T. P. Tang.
``New Algorithms for Improved Transcendental Functions on IA-64.''
\emph{14th IEEE Symposium on Computer Arithmetic}, pp.~4--11, 1999.
\href{https://doi.org/10.1109/ARITH.1999.762822}{Publisher record}.
\end{thebibliography}

\clearpage
\appendix
\section{Literal constant specification}
\label{app:constants}

Each row below defines $C_i=(-1)^s S2^q$, where $S$ is the listed
hexadecimal integer and $q$ is the binary scale. Rows 114--117 form the
short polynomial; rows 118--123 form the long polynomial. Rows 125--156
are the anchors $\operatorname{atan}(n/32)$ at $i=124+n$ as represented in
the P5 ROM, not values recomputed with a host math library. The table is
generated from the checked C literals and cross-checked against the public
transcription; Section~\ref{sec:rom} gives the exponent conversion.

% Generated from the literal C entries, checked against the public TSV.
{\footnotesize
\begin{longtable}{@{}rrrl@{}}
\toprule
ROM row $i$ & Sign $s$ & Scale $q$ & Integer significand $S$ (hex) \\
\midrule
\endfirsthead
\toprule
ROM row $i$ & Sign $s$ & Scale $q$ & Integer significand $S$ (hex) \\
\midrule
\endhead
19 & 0 & -65 & \texttt{6487ed5110b4611a6} \\
20 & 0 & -66 & \texttt{6487ed5110b4611a6} \\
\midrule
114 & 1 & -68 & \texttt{555555555555535f0} \\
115 & 0 & -69 & \texttt{6666666664208b016} \\
116 & 1 & -69 & \texttt{492491e0653ac37b8} \\
117 & 0 & -70 & \texttt{71b83f4133889b2f0} \\
\midrule
118 & 1 & -68 & \texttt{55555555555555543} \\
119 & 0 & -69 & \texttt{66666666666616b73} \\
120 & 1 & -69 & \texttt{4924924920fca4493} \\
121 & 0 & -70 & \texttt{71c71c4be6f662c91} \\
122 & 1 & -70 & \texttt{5d16e0bde0b12eee8} \\
123 & 0 & -70 & \texttt{4e403be3e3c725aa0} \\
\midrule
125 & 0 & -72 & \texttt{7ff556eea5d892a14} \\
126 & 0 & -71 & \texttt{7fd56edcb3f7a71b6} \\
127 & 0 & -70 & \texttt{5fb860980bc43a305} \\
128 & 0 & -70 & \texttt{7f56ea6ab0bdb7196} \\
129 & 0 & -69 & \texttt{4f5bbba31989b161a} \\
130 & 0 & -69 & \texttt{5ee5ed2f396c089a4} \\
131 & 0 & -69 & \texttt{6e435d4a498288118} \\
132 & 0 & -69 & \texttt{7d6dd7e4b203758ab} \\
133 & 0 & -68 & \texttt{462fd68c2fc5e0986} \\
134 & 0 & -68 & \texttt{4d89dcdc1faf2f34e} \\
135 & 0 & -68 & \texttt{54c2b6654735276d5} \\
136 & 0 & -68 & \texttt{5bd86507937bc239c} \\
137 & 0 & -68 & \texttt{62c934e5286c95b6d} \\
138 & 0 & -68 & \texttt{6993bb0f308ff2db2} \\
139 & 0 & -68 & \texttt{7036d3253b27be33e} \\
140 & 0 & -68 & \texttt{76b19c1586ed3da2b} \\
141 & 0 & -68 & \texttt{7d03742d50505f2e3} \\
142 & 0 & -67 & \texttt{4195fa536cc33f152} \\
143 & 0 & -67 & \texttt{4495766fef4aa3da8} \\
144 & 0 & -67 & \texttt{47802eaf7bfacfcdb} \\
145 & 0 & -67 & \texttt{4a563964c238c37b1} \\
146 & 0 & -67 & \texttt{4d17c07338deed102} \\
147 & 0 & -67 & \texttt{4fc4fee27a5bd0f68} \\
148 & 0 & -67 & \texttt{525e3e8c9a7b84921} \\
149 & 0 & -67 & \texttt{54e3d5ee24187ae45} \\
150 & 0 & -67 & \texttt{5756261c5a6c60401} \\
151 & 0 & -67 & \texttt{59b598e48f821b48b} \\
152 & 0 & -67 & \texttt{5c029f15e118cf39e} \\
153 & 0 & -67 & \texttt{5e3daef574c579407} \\
154 & 0 & -67 & \texttt{606742dc562933204} \\
155 & 0 & -67 & \texttt{627fd7fd5fc7deaa4} \\
156 & 0 & -67 & \texttt{6487ed5110b4611a6} \\
\bottomrule
\end{longtable}
}


\clearpage
\section{Programmer's pseudocode}
\label{app:pseudocode}

The following blocks are copied mechanically from the standalone
\path{PSEUDOCODE.md}. They are executable reference pseudocode with exact
\texttt{Fraction} arithmetic; the documentation verifier binds
\texttt{ROM} to Appendix~\ref{app:constants}. No production numerical
function is called by this reference. Internal rounding has an unbounded
exponent range; only \texttt{pack\_angle} applies raw80 output spacing.
The returned tuple is (raw result, C1, exception flags, pre-load flags).
IE=1, DE=2, ZE=4, OE=8, UE=16 and PE=32. The numerical wrapper assumes the
clear, masked, valid-stack capture contract, not arbitrary prior FPU state.

% Generated verbatim from PSEUDOCODE.md; do not hand-edit.
\subsection{Raw80 values, classification and exact arithmetic}
\par\noindent\begin{minipage}{\linewidth}
\lstinputlisting[style=pseudocode,firstline=1,lastline=20]{generated/pseudocode-1.py}
\end{minipage}\par
\par\noindent\begin{minipage}{\linewidth}
\lstinputlisting[style=pseudocode,firstline=21,lastline=34]{generated/pseudocode-1.py}
\end{minipage}\par
\subsection{Internal rounding and final raw80 encoding}
\par\noindent\begin{minipage}{\linewidth}
\lstinputlisting[style=pseudocode,firstline=1,lastline=24]{generated/pseudocode-2.py}
\end{minipage}\par
\par\noindent\begin{minipage}{\linewidth}
\lstinputlisting[style=pseudocode,firstline=25,lastline=43]{generated/pseudocode-2.py}
\end{minipage}\par
\subsection{Two interleaved polynomial chains}
\par\noindent\begin{minipage}{\linewidth}
\lstinputlisting[style=pseudocode,firstline=1,lastline=21]{generated/pseudocode-3.py}
\end{minipage}\par
\subsection{Finite reduction, dispatch and quadrant restoration}
\par\noindent\begin{minipage}{\linewidth}
\lstinputlisting[style=pseudocode,firstline=1,lastline=27]{generated/pseudocode-4.py}
\end{minipage}\par
\subsection{Masked architectural wrapper}
\par\noindent\begin{minipage}{\linewidth}
\lstinputlisting[style=pseudocode,firstline=1,lastline=42]{generated/pseudocode-5.py}
\end{minipage}\par


\end{document}
